\chapter{Methodology} 
\label{chap:methodology} 

This chapter describes the methodological approach used to develop and evaluate the proposed security testing framework. Building upon the theoretical foundations presented in Chapter \ref{chap:fundamentals}, the research follows the Design Science Research Methodology (DSRM) proposed by \citeauthor{Peffers} DSRM provides a structured process for addressing practical problems through the design, development, demonstration, and evaluation of artifacts \cite[p. 46]{Peffers}. The objective of this research is the development of a standards-based functional security testing framework for industrial gateways. The framework is intended to enable the traceable, repeatable, and systematic verification of security functions derived from the requirements of Regulation (EU) 2023/1230, prEN 50742, and the IEC 62443 series. Within the scope of this thesis, the activities of the DSRM process are mapped to the corresponding chapters. Problem identification and objective definition are presented in Chapter \ref{chap:introduction}. The design and development phase is addressed in this chapter through the derivation of security requirements and the development of the proposed testing framework. The demonstration and evaluation phases are conducted in Chapter \ref{chap:implementation_results_evaluation}, where the framework is applied to an industrial gateway as a case study. The communication phase is represented by the documentation, discussion, and dissemination of the research results through this thesis.

\section{Derivation of Security Requirements}
\label{sec:security_requirements}

The aim is to produce a structured, traceable set of product requirements for an industrial gateway at the IT/OT interface, grounded in the legal obligations of Regulation~(EU) 2023/1230. The complete requirement catalogue, prEN~50742 mapping, and IEC~62443-4-2 mapping are documented in Appendix~\ref{app:requirement_derivation}.


The derivation follows a three-stage hierarchical chain: 

\begin{enumerate}
	\item \hyperref[subsec:eu_machinery_regulation]{\textbf{Regulation~(EU) 2023/1230}} defines legally binding protection objectives for machinery, including explicit cybersecurity obligations.
	\item \hyperref[subsec:pren_50742]{\textbf{prEN~50742}} operationalizes these objectives for connected machinery and links them to the IEC~62443 series.
	\item \hyperref[subsec:iec_62443]{\textbf{IEC~62443-4-2}} specifies component-level technical security requirements (CRs, REs, EDRs) that can be verified by functional testing.
\end{enumerate}

None of the three sources is independently sufficient. Regulation~(EU) 2023/1230 defines mandatory objectives but not component-level technical specifications. IEC~62443-4-2 provides precise technical controls but does not by itself establish which are required for Machinery Regulation conformity. PrEN~50742 preserves legal traceability while introducing technically interpretable obligations that can be mapped to IEC~62443-4-2.


\hyperref[subsec:eu_machinery_regulation]{\textbf{Regulation~(EU) 2023/1230}} establishes that machinery certified under a cybersecurity certification scheme pursuant to Regulation~(EU) 2019/881 is presumed to comply with the essential health and safety requirements of Annex~III insofar as these requirements are covered by the certification. The cybersecurity provisions explicitly addressed by this presumption are:

\begin{itemize}
	\item \textbf{Section~1.1.9} (Protection against corruption): obligations concerning safety-relevant hardware interfaces, software and data protection, and evidence collection.
	\item \textbf{Section~1.2.1} (Safety and reliability of control systems): obligations concerning resilience against malicious third-party attempts and logging of safety-relevant interventions.
\end{itemize}



Each regulatory clause was decomposed analytically into candidate requirements. The decomposition applied the following four criteria:

\begin{enumerate}
	\item \textbf{Traceability:} The requirement must trace directly to an explicit obligation in \hyperref[subsec:eu_machinery_regulation]{\textbf{Regulation~(EU) 2023/1230}}.
	\item \textbf{Atomicity:} Each requirement must express exactly one verifiable technical property. Clauses combining multiple independent obligations were split into separate requirements.
	\item \textbf{Verifiability:} The requirement must be testable in principle by a functional test or a clearly delimited technical assessment.
	\item \textbf{Product relevance:} The requirement must describe a property of the gateway as a product, not an organizational, documentation, or purely functional-safety obligation.
\end{enumerate}



Each candidate requirement was compared with \hyperref[subsec:pren_50742]{\textbf{prEN~50742}} to determine whether the draft standard provides a technically defensible interpretation for the machinery domain. Four mapping types were defined: 
\begin{itemize}
	\item \textbf{Direct}: clause-level match
	\item \textbf{Partial}: scope or trigger condition differs 
	\item \textbf{Indirect}: related principle, no clause-level match
	\item \textbf{Not identified}: confirmed gap or conflict
\end{itemize}


Where \hyperref[subsec:pren_50742]{\textbf{prEN~50742}} provided clause-level support, the corresponding IEC~62443-4-2 CRs, REs, and EDRs were identified. One additional mapping type was introduced for IEC~62443-4-2: \textbf{No direct mapping} (no CR addresses the obligation).

Controls were selected only where the connection to the regulatory source could be substantiated through the derivation chain. The role of the Foundational Requirements (see Section~\ref{subsec:iec_62443}) was classificatory. The selection of specific CRs was driven by the protection goals established in the earlier stages.

Confirmed normative gaps arise for two requirements within the scope of Sections~1.1.9 and~1.2.1. RQ-010 addresses resilience against malicious third-party attempts at a system level. IEC~62443-4-2 assigns such process-level obligations to IEC~62443-4-1, and no component-level CR provides a direct mapping. RQ-013 presents a confirmed conflict: the MVO requires logging of uploaded safety software versions, whereas prEN~50742 explicitly excludes binary software versions from the tracing log. These two gaps define the positions where the standard-based test approach must be supplemented by product-specific verification.

The detailed mapping tables are provided in Appendix~\ref{app:sec:mapping_pren50742} (prEN~50742) and Appendix~\ref{app:sec:mapping_iec62443} (IEC~62443-4-2).


The derivation produced 14 requirements (RQ-001 to RQ-014). Each requirement was classified by the security objective it primarily expresses, corresponding to the security goals and Foundational Requirements introduced in Sections~\ref{subsec:security_goals} and~\ref{subsec:iec_62443}. Table~\ref{tab:req_categories} summarizes the resulting categories and the IEC~62443-4-2 coverage achieved.

\renewcommand{\arraystretch}{1.2}
\begin{table}[htbp]
	\caption{Summary of Derived Security Requirement Categories}
	\label{tab:req_categories}
	\centering
	\begin{tabularx}{\textwidth}{|l | X | X | X|}
		\toprule
		\textbf{Category} & \textbf{FR} & \textbf{Req.\ IDs} & \textbf{IEC~62443-4-2 Coverage} \\
		\midrule
		Secure Communication        & FR3              & RQ-001                                   & Partial \\
		Data and Software Integrity & FR3              & RQ-002, RQ-005                           & Direct (RQ-005); Partial (RQ-002) \\
		Configuration Protection    & FR1, FR2, FR7    & RQ-004, RQ-006, RQ-007, RQ-011           & Partial \\
		Logging and Auditability    & FR2, FR3         & RQ-003, RQ-008, RQ-009, RQ-012, RQ-013  & Direct (RQ-003, RQ-008, RQ-009, RQ-012); None (RQ-013) \\
		Access Control              & FR3, FR6         & RQ-014                                   & Partial \\
		Availability / Resilience   & -              & RQ-010                                   & None \\
		\bottomrule
	\end{tabularx}
\end{table}

Of the 14 requirements, 5 achieve a Direct IEC~62443-4-2 mapping, 7 achieve a Partial mapping, and 2 obtain no direct IEC~62443-4-2 counterpart. The complete requirement catalogue with source text, security objective, requirement type, and rationale is provided in Appendix~\ref{app:sec:mvo_requirements}.


The derived requirement base defines the test scope of this thesis. Requirements with a Direct or Partial IEC~62443-4-2 mapping form the primary basis for functional test case design, since they correspond to established, technology-neutral component security capabilities. Partially mapped requirements require interpretive judgment when converted into executable test procedures.

Requirements classified as \emph{No direct mapping} define the normative boundary of the standard-based approach. At these positions, the test methodology must incorporate product-specific verification procedures or explicitly document the remaining conformity gap.

The categorization in Table~\ref{tab:req_categories} structures the subsequent test case design. Each security function category is associated with a set of test strategies drawn from the testing approaches described in Section~\ref{sec:security_testing_approaches}. Requirements within a category with identical IEC~62443-4-2 anchors share the same verification logic and can be addressed by a common test procedure.



\section{Functional Security Testing Framework Architecture}
\label{sec:framework_architecture}

\subsection{Scope and Boundaries}
\label{subsec:scope_boundaries}

This thesis is delimited to \textbf{Functional verification}, i.e., the testing of implemented security functions against the requirement baseline (Section~\ref{sec:security_requirements}), as opposed to \textbf{robustness testing} (fuzzing, penetration testing) at the gateway's IT/OT interfaces, as already distinguished in Section~\ref{sec:security_testing_approaches} and Section~\ref{subsec:vulnerabilities_fuzzing_pentesting}.

Within this scope, the framework verifies that the security functions and countermeasures mandated by RQ-001 to RQ-014 are actually implemented and operate as specified at the interfaces identified as gateway assets in Section~\ref{subsec:attack_surface_asset_identification}. Fuzzing, penetration testing, and open-ended vulnerability scanning are explicitly out of scope, consistent with the exclusions already defined in Chapter~\ref{chap:introduction}. These techniques search for unknown weaknesses rather than verifying a specified, requirement-bound behavior, and therefore answer a different question than the one posed by this framework \cite[p.8, pp. 29]{richier2011security}.


\subsection{Traceability Model}
\label{subsec:traceability_model}

The traceability chain reflects the perspective of the tester executing the framework, linking an abstract regulatory obligation to a concrete, evidenced test result. Each relationship in the chain is one-to-many. One requirement may apply to several assets, and one test category may be instantiated by several test cases.

Requirement $\rightarrow$ Asset $\rightarrow$ Threat $\rightarrow$ Test Category $\rightarrow$ Test Method $\rightarrow$ Test Case $\rightarrow$ Evidence and Verdict

The chain is read as follows. 
The \textbf{MVO/IEC Requirement} is one of the RQ-001 to RQ-014 requirements together with its mapped IEC 62443-4-2 CR/RE/EDR (Section~\ref{sec:security_requirements}).
 
The \textbf{Asset} is the concrete gateway asset (Section~\ref{subsec:test_object}) to which the requirement applies, anchoring the abstract obligation to a testable component of the test object. 

The \textbf{STRIDE Threat} is the threat category (Section~\ref{subsec:stride}) that the requirement is intended to counter at that asset. 

The \textbf{Test Category} groups requirement-asset-threat combinations by verification logic, as defined in Section~\ref{subsec:test_case_taxonomy}. 

The \textbf{Test Method} is the established functional testing technique assigned to that category (Section \ref{sec:security_testing_approaches}), so that the tester verifies the correct operation of a planned security mechanism rather than searching for unknown flaws \cite[p. 82]{1341418}. 

The \textbf{Test Case} is the concrete instantiation of that method, applied with a specific tool and procedure to the test object (Chapter~\ref{chap:implementation_results_evaluation}). 

Finally, \textbf{Evidence and Verdict} comprises the recorded artifacts (logs, captures, screenshots) and the discrete verdict assigned to the test case (Section \ref{subsec:verdict_scheme}), in line with the documentation and reporting discipline required of structured security testing \cite[p. 7-3, p. 7-5]{nist_sp800_115}.

For the design phase, an analogous forward traceability from threat to control is already established by the TARA workflow (Section \ref{subsec:tara}), which proceeds from threat and impact analysis to a documented risk treatment decision \cite[pp. 46]{doCarmo2025}. 
The chain defined here starts from the mitigation already established through the TARA process and the requirement derivation described in Section~\ref{sec:security_requirements} (see Section~\ref{subsec:risk_treatment_strategies}). It then carries this mitigation forward, from the tester's perspective, through the concrete asset, threat, and test method to an executed and evidenced verdict. Consequently, no mitigation is credited as effective solely on the basis of its specification.

Every test case therefore carries an unbroken reference back to its originating regulatory clause. This traceability chain is the structural precondition for the coverage metric introduced in Section~\ref{subsec:coverage_metric} and for the technical documentation obligation under Article 10(2) and Annex IV Part A of Regulation (EU) 2023/1230 \cite{EU2023_1230}.
